\section{Discussion and Limitations}

The benchmark supports a narrow but coherent claim: once adaptation has a non-negligible energy cost, the best adaptation schedule can depend on the deployment regime. That is the most useful scientific signal in the sample. It also aligns with the review pressure from the earlier workflow stages, which repeatedly pushed the project away from ``invent a new TTA method'' and toward ``show a regime where adaptation scheduling matters.'' The right reading of the result is therefore not ``battery-gated adaptation is the new default,'' but ``battery-aware scheduling is a variable the deployment story must explain.''

At the same time, the limitations are substantial. The benchmark is synthetic. The numbers were created to continue the workflow honestly inside the repository, not to stand in for a real quadruped evaluation. The controller schematic and regime map remain placeholders because the repository does not contain an actual adaptation implementation to visualize. A real project would need at least three follow-on steps: instrument a real or simulated locomotion stack, validate the same crossover across multiple disturbance families, and replace the placeholder figures with real plots and system diagrams. Until then, the paper should be read as a workflow artifact with a plausible research hypothesis, not as scientific evidence about quadruped deployment.

These limitations are not incidental. They are the reason the paper is framed as an end-to-end workflow demonstration rather than a publishable robotics claim. The point of the sample is to prove artifact completeness and to make the final paper stage concrete in the repository itself.
